\documentclass[12pt]{article}
\usepackage[utf8]{inputenc}
\usepackage[T1]{fontenc}
\usepackage{hyperref}
\title{Penrhyn Language Resources for K-12 Curriculum Development}
\author{Compiled by Claude Sonnet 5.5 (claude-sonnet-5-5) for Andr\'as Kornai}
\date{}
\begin{document}
\maketitle

\section{Introduction}
Penrhyn (Tongareva, ISO 639-3 pnh) is an endangered Polynesian language of the northern Cook Islands, with about 200 speakers reported in the 2011 census. Its classification is contested: some sources treat it as a dialect of Cook Islands M\=aori, others as a separate language. Published material is very limited, and this report lists only what the searches turned up.

\section{Primary Educational Resources}
No Penrhyn primers, textbooks or school curricula were found in the searches. The online dictionary below is the closest to a classroom-usable resource.
\begin{itemize}
\item Dictionary of Cook Islands Languages, Penrhyn section (searchable in both directions): \url{https://penrhyn.cookislandsdictionary.com/}
\item Cook Islands Dictionary mobile app (AUT University), which includes a Tongarevan (Penrhyn) dictionary: \url{https://apps.apple.com/us/app/id6754361002}
\end{itemize}

\section{Linguistic and Reference Materials}
\begin{itemize}
\item The online dictionary is based on Norio Shibata's 2003 \emph{Penrhyn-English Dictionary} (out of print); see the about page: \url{https://penrhyn.cookislandsdictionary.com/about}
\item Further resources page of the dictionary site: \url{https://penrhyn.cookislandsdictionary.com/resources}
\item Yasuda (1968), \emph{The structure of the Penrhyn phrase}, University of Hawaii MA thesis (typescript, restricted reproduction): \url{https://scholarspace.manoa.hawaii.edu/items/d7ff038e-ae5d-42bc-93ee-6e79d8457a3b}
\item Overview article: \url{https://en.wikipedia.org/wiki/Penrhyn_language}
\item Writing system overview: \url{https://mail.omniglot.com/writing/penrhyn.htm}
\end{itemize}

\section{Religious and Community Resources}
No Penrhyn scripture translations or community organisation pages were found in the searches.

\section{Digital and Audio Resources}
\begin{itemize}
\item Endangered Languages Project entry: \url{https://www.endangeredlanguages.com/lang/2994}
\item Alternative entry: \url{https://endangeredlanguages.com/node/97004}
\end{itemize}
No audio or video collections were found in the searches.

\section{Curriculum Development Context}
The language is not known to be taught in schools, and the speaker base is small and mainly adult. A curriculum would have to be built largely from the Shibata dictionary material and from new documentation with speakers. Buck's \emph{Ethnology of Tongareva} (Bishop Museum Bulletin 92, 1932) is recommended in a search summary as cultural background, but no URL was seen for it.

\end{document}